\chapter{Metodología}

\section{Enfoque de desarrollo}

La metodología combina desarrollo incremental de software y validación técnica del flujo. El objetivo no es entrenar un modelo financiero ni establecer una comparativa de servicios LLM, sino construir un prototipo capaz de transformar consultas financieras históricas en cálculos reproducibles y explicaciones prudentes.

La unidad de análisis es el flujo completo: entrada validada, planificación, generación de código, validación de seguridad, ejecución controlada e interpretación final. Esta decisión permite observar no sólo si una respuesta parece correcta, sino en qué punto del sistema se produce un acierto o un fallo.

\section{Proceso incremental}
\label{sec:evolucion_mvp}

El desarrollo se organizó en iteraciones:

\begin{enumerate}
    \item Definir una entrada estructurada con consulta, tickers, fechas, intervalo y rutas a CSV locales.
    \item Validar campos mínimos antes de invocar el modelo.
    \item Separar el uso del LLM en tres agentes: planificación, generación de código e interpretación.
    \item Exigir contratos JSON entre agentes y entre el código generado y el sistema.
    \item Validar el código antes de ejecutarlo, limitando importaciones y operaciones permitidas.
    \item Ejecutar el script sobre datos históricos locales y capturar salida estándar, errores y artefactos.
    \item Compactar salidas largas antes de la interpretación final cuando sea necesario.
    \item Revisar los resultados mediante métricas técnicas y, como línea futura, métricas subjetivas más reforzadas.
\end{enumerate}

Este enfoque separa tres planos: datos, ejecución e interpretación. El plano de datos asegura que la consulta se resuelve sobre ficheros conocidos. El plano de ejecución controla el código generado. El plano interpretativo convierte resultados estructurados en una respuesta legible.

\section{Criterios de validación}

La validación del prototipo se organiza como una cadena de controles. Cada control detecta un tipo de fallo antes de que afecte a la respuesta final.

\begin{table}[H]
\centering
\small
\begin{tabularx}{\textwidth}{p{0.22\textwidth} X X}
\toprule
\textbf{Nivel de validación} & \textbf{Qué comprueba} & \textbf{Riesgo que reduce} \\
\midrule
Entrada & Consulta no vacía, tickers, fechas, intervalo y CSV disponibles. & Evita lanzar el flujo con peticiones incompletas o imposibles de resolver. \\
Plan analítico & Nivel A/B/C, objetivo, métricas, columnas y salidas esperadas. & Reduce desalineamientos entre lo pedido y lo calculado. \\
Código generado & Función \texttt{main()}, lectura de \texttt{argv[1]}, importaciones permitidas y salida JSON. & Reduce riesgos de seguridad y salidas no parseables. \\
Ejecución & Código ejecutado, errores, métricas calculadas y salida estructurada. & Distingue fallo de generación, fallo de ejecución y resultado válido. \\
Interpretación & Uso exclusivo de resultados calculados, claridad, limitaciones y prudencia financiera. & Reduce alucinaciones y recomendaciones no permitidas. \\
Trazabilidad & Conservación de entrada, plan, código, salida y errores. & Permite auditar y repetir una ejecución. \\
\bottomrule
\end{tabularx}
\caption{Criterios de validación aplicados al flujo del prototipo.}
\label{tab:criterios_validacion_flujo}
\end{table}

\section{Niveles de dificultad}

La batería de evaluación se organiza en tres niveles:

\begin{itemize}
    \item \textbf{Nivel A}: consultas directas, sin formato especial. Se espera una selección mínima de métricas esenciales.
    \item \textbf{Nivel B}: consultas que piden más explicación, comparación, tabla o una salida visual principal. Se espera una salida estructurada más rica.
    \item \textbf{Nivel C}: consultas profesionales, multicriterio o con formato impuesto. Se espera coordinación de varias métricas, tablas, datos visuales si se solicitan y limitaciones.
\end{itemize}

La batería se equilibra con cinco consultas por nivel. Dentro de cada nivel, los ejemplos ascienden en complejidad: primero casos directos, después comparaciones, luego riesgo o ventanas temporales más exigentes, y finalmente consultas con formato más rígido.

\section{Métricas de evaluación}

La evaluación actual se centra en métricas técnicas:

\begin{itemize}
    \item estado final del workflow;
    \item llegada o no a planificación, generación de código, ejecución e interpretación;
    \item errores de configuración, formato, validación, ejecución o interpretación;
    \item presencia de salida estructurada;
    \item ausencia de recomendaciones de inversión;
    \item tiempo de ejecución por caso.
\end{itemize}

Las métricas subjetivas quedan planteadas como trabajo futuro. Se consideran relevantes para valorar claridad, utilidad percibida, adecuación al nivel A/B/C y calidad comunicativa, pero requieren una rúbrica más sólida antes de usarse como evidencia principal. Por ello, en esta memoria no se presentan como resultados concluyentes.

\section{Revisión humana}

La revisión humana no sustituye a la validación técnica. Su papel es detectar si una respuesta técnicamente completa resulta insuficiente, ambigua o poco útil. En futuras iteraciones se reforzará mediante una rúbrica explícita y repetible, separando comprensión de la consulta, selección de métricas, fidelidad a la salida ejecutada y claridad de la interpretación.
